\documentclass[12pt,a4paper]{article}
\usepackage[utf8]{inputenc}
\usepackage[T1]{fontenc}
\usepackage{amsmath,amssymb,amsthm}
\usepackage{geometry}
\usepackage{titlesec}
\usepackage{fancyhdr}
\usepackage{hyperref}
\usepackage{etoolbox}
\usepackage{array}

\geometry{margin=1in}

\pagestyle{fancy}
\fancyhf{}
\fancyhead[C]{\small Cayley Mathematical Papers --- Vol.~III, pp.~201--250}
\fancyfoot[C]{\thepage}
\renewcommand{\headrulewidth}{0.4pt}

\theoremstyle{definition}
\newtheorem*{remark}{Remark}

\titleformat{\section}{\large\bfseries}{\thesection.}{0.5em}{}
\titleformat{\subsection}{\normalsize\bfseries}{\thesubsection.}{0.5em}{}

\begin{document}

\begin{center}
{\Large\bfseries The Collected Mathematical Papers of\\[0.5em] Arthur Cayley, Vol.~III}\\[2em]
{\large Pages 201--250}\\[1em]
{\normalsize Transcribed from the Cambridge University Press edition}
\end{center}

\vspace{2em}

\tableofcontents
\newpage


\section*{195. Report on the Recent Progress of Theoretical Dynamics (concluded)}
\addcontentsline{toc}{section}{195. Report on the Recent Progress of Theoretical Dynamics (concluded)}

\noindent\textit{(Conclusion from previous pages.)}

\medskip

I remark in conclusion, that the differential equations of dynamics (including in
the expression, as I have done throughout the report, the generalized Lagrangian and
Hamiltonian forms) constitute really one of the chief classes of differential equations
occupying the attention of geometers. The greater part of what has been done with
respect to the general theory of a system of differential equations is due to Jacobi,
and he has also considered in particular, besides the differential equations of dynamics,
the Pfaffian system of differential equations (including therein the system of differential
equations which arise from any partial differential equation of the first order), and
the so-called isoperimetric system of differential equations, that is, the system arising
from any problem in the calculus of variations. In a systematic treatise it would be
proper to commence with the general theory of a system of differential equations, and
as a branch of this general theory, to consider the generalized Hamiltonian system;
and in relation thereto to develope the various theorems which have a dynamical
application. It would be shewn that the generalized Lagrangian form would be transformed
into the Hamiltonian form, but the first-mentioned form would, I think,
properly be treated as a particular case of the isoperimetric system of differential
equations.

\bigskip

\begin{center}
\textit{List of Memoirs and Works above referred to.}
\end{center}

\medskip

\noindent\textsc{Lagrange}. M\'ecanique Analytique. 1st edition. 1788.

\smallskip

\noindent\textsc{Laplace}. M\'ecanique C\'eleste, t.~i. 1799.

\smallskip

\noindent\textsc{Poisson}. Sur les in\'egalit\'es s\'eculaires des moyens mouvemens des plan\`etes. Read to
the Institute 20th June, 1808.---\textit{Journ. \'Ecole Polyt.} t.~VIII. pp.~1--56. 1809.

\smallskip

\noindent\textsc{Laplace}. M\'emoire.... Read to the Bureau of Longitudes 17th Aug., 1808.---Forms
an Appendix to the 3rd volume of the M\'ecanique C\'eleste. 1809.

\smallskip

\noindent\textsc{Lagrange}. M\'emoire sur la th\'eorie des variations des \'el\'emens des plan\`etes et en
particulier des variations des grands axes de leurs orbites. Read to the Bureau
of Longitudes the 17th Aug. and to the Institute the 22nd Aug. 1808.---M\'em.
de l'Instit. 1808, pp.~1--72. 1809.

\smallskip

\noindent\textsc{Lagrange}. M\'emoire sur la th\'eorie g\'en\'erale de la variation des constantes arbitraires
dans tous les probl\`emes de la M\'ecanique. Read to the Institute 13th March,
1809.---M\'em. de l'Instit, 1808, pp.~257--302 (includes an undated addition), and
there is a Supplement also without date, pp.~363, 364. 1809.

\smallskip

\noindent\textsc{Poisson}. M\'emoire sur la variation des constantes arbitraires dans les questions de
M\'ecanique. Read to the Institute 16th Oct., 1809.---Journ. \'Ecole Polyt., t.~VIII.
pp.~266--344. 1809.

\smallskip

\noindent\textsc{Lagrange}. Seconde m\'emoire sur la variation des constantes arbitraires dans les
probl\`emes de M\'ecanique, dans lequel on simplifie l'application des formules g\'en\'erales
\`a ces probl\`emes. Read to the Institute 19th Feb., 1810.---M\'em. de l'Instit. 1809,
pp.~343--352. 1810.

\smallskip

\noindent\textsc{Lagrange}. M\'ecanique Analytique. 2nd edition, t.~i. 1811; t.~II. 1815.

\smallskip

\noindent\textsc{Poisson}. M\'emoire sur la variation des constantes arbitraires dans les questions de
M\'ecanique. Read to the Academy, 2nd Sept., 1816.---Bulletin de la Soc. Philom.,
1816, p.~109; and also, M\'em. de l'Instit., t.~I. pp.~1--70. 1816.

\smallskip

\noindent\textsc{Hamilton, Sir W. R.} On a general method in Dynamics, by which the study of all
free systems of attracting or repelling points is reduced to the search and differentiation
of one central relation or characteristic function.---Phil. Trans. for 1834,
pp.~247--308. 1834.

\smallskip

\noindent\textsc{Hamilton, Sir W. R.} Second Essay on a general method in Dynamics.---Phil. Trans.
for 1835, pp.~95--144. 1835.

\smallskip

\noindent\textsc{Poisson}. Remarques sur l'Int\'egration des \'equations diff\'erentielles de la dynamique.---
Liouville, t.~II. pp.~317--337. 1837.

\smallskip

\noindent\textsc{Jacobi}. \"Uber die Reduction der Integration der partiellen Differentialgleichungen
erster Ordnung zwischen irgend einer Zahl Variabeln auf die Integration eines
einzigen Systemes gew\"ohnlicher Differentialgleichungen.---Crelle, t.~XVII. pp.~97--162,
and translated into French, Liouville, t.~III. pp.~60--96, and 161--201. 1837.

\smallskip

\noindent\textsc{Jacobi}. Note sur l'int\'egration des \'equations diff\'erentielles de la dynamique.---Comptes
Rendus, t.~v. p.~61. 1837.

\smallskip

\noindent\textsc{Jacobi}. Nova Theorema der analytischen Mechanik.---Monatsbericht of the Academy
of Berlin for 1838 (paper is dated 21st Nov., 1836); and Crelle, t.~XXX. pp.~117--120.
1838.

\smallskip

\noindent\textsc{Jacobi}. Lettre adress\'ee \`a M. le Pr\'esident de l'Acad\'emie des Sciences \`a Paris.---
Comptes Rendus, t.~XI. p.~529; Liouville, t.~V. pp.~350--351 (with an addition by
Liouville, pp.~351--355). 1840.

\smallskip

\noindent\textsc{Binet}. M\'emoires sur la variation des constantes arbitraires dans les formules g\'en\'erales
de la dynamique et dans les syst\`emes d'\'equations analogues plus \'etendus.---Journ.
\'Ecole Polyt., t.~XVII. pp.~1--94. 1841.

\smallskip

\noindent\textsc{Jacobi}. Sur un nouveau principe g\'en\'eral de la M\'ecanique Analytique.---Comptes
Rendus, t.~XV. pp.~202--205. 1842.

\smallskip

\noindent\textsc{Jacobi}. De motu puncti singularis.---Crelle, t.~XXIV. pp.~5--27. 1842.

\smallskip

\noindent\textsc{Maurice}. M\'emoire sur la variation des constantes arbitraires nous l'\'etat \'etablie dans
sa g\'en\'eralit\'e les coh\'erents de Lagrange et celui de Poisson. Read to the Academy
the 1st June, 1844.---M\'em. de l'Instit., t.~XIX. pp.~553--628. 1844.

\smallskip

\noindent\textsc{Jacobi}. Theoria novi multiplicatoris systemati \'equations differentialium vulgarium
applicandi.---Crelle, t.~XXVII. pp.~199--268; and t.~XXIX. pp.~213--279, and 333--376.
1844.

\smallskip

\noindent\textsc{Dedeyen}. D\'emonstration de deux th\'eor\`emes de M. Jacobi, application au probl\`eme des
perturbations plan\'etaires. Thesis presented to the Faculty of Sciences the 3rd
April, 1848.---Liouville, t.~XIII. pp.~207--411. 1848.

\smallskip

\noindent\textsc{Serret}. Sur l'int\'egration des \'equations diff\'erentielles du mouvement d'un point mat\'eriel.---
Comptes Rendus, t.~XXVI. pp.~605--608. 1848.

\smallskip

\noindent\textsc{Serret}. Sur l'int\'egration des \'equations diff\'erentielles de la dynamique.---Comptes Rendus,
t.~XXVI. pp.~639--641. 1848.

\smallskip

\noindent\textsc{Sturm}. Note sur l'int\'egration des \'equations g\'en\'erales de la dynamique.---Comptes
Rendus, t.~XXVI. pp.~658--673. 1848.

\smallskip

\noindent\textsc{Ostrogradsky}. Sur les int\'egrales des \'equations g\'en\'erales de la dynamique.---M\'elanges
de l'Acad. de St. P\'etersbourg, t.~II. Oct. 1848.

\smallskip

\noindent\textsc{[Ostrogradsky]}. M\'emoire sur les \'equations diff\'erentielles relatives au probl\`eme des
isop\'erim\`etres.---M\'em. de l'Acad. de St. P\'etersb., t.~VI. pp.~385--517. 1850.]

\smallskip

\noindent\textsc{Brassinne}. Th\'eor\`emes relatifs \`a une classe d'\'equations diff\'erentielles simultan\'ees analogues
\`a un th\'eor\`eme employ\'e par Lagrange dans la th\'eorie des perturbations.---Liouville,
t.~XVI. pp.~382--388. 1851.

\smallskip

\noindent\textsc{Bertrand}. M\'emoire sur les int\'egrales communes \`a plusieurs probl\`emes de M\'ecanique.
Presented to the Academy the 12th May, 1851.---Liouville, t.~XVII. pp.~121--174.
1852.

\smallskip

\noindent\textsc{Bertrand}. M\'emoire sur l'int\'egration des \'equations diff\'erentielles de la dynamique.---
Liouville, t.~XVII. pp.~393--436. 1852.

\smallskip

\noindent\textsc{Bertrand}. Sur un nouveau th\'eor\`eme de M\'ecanique Analytique.---Comptes Rendus,
t.~XXXV. pp.~698--699. 1852.

\smallskip

\noindent\textsc{Bertrand}'s notes VI. and VII. to the third edition of the M\'ecanique Analytique,
t.~I. pp.~423--428, t.~II. note VI.---Sur les \'equations diff\'erentielles des probl\`emes de
M\'ecanique et la forme que l'on peut donner \`a leurs int\'egrales; and note VII.---
Sur un th\'eor\`eme de Poisson. 1853.

\smallskip

\noindent\textsc{Briosch\'\i}. Sulla variazione delle costanti arbitrarie nei problemi della Dinamica.---
Tortolini, Annali, t.~IV. pp.~298--311. 1853.

\smallskip

\noindent\textsc{Briosch\'\i}. Intorno ad un teorema di Meccanica.---Tortolini, Annali, t.~IV. pp.~351--354.
1853.

\smallskip

\noindent\textsc{Liouville}. Note sur l'int\'egration des \'equations diff\'erentielles de la dynamique, pr\'esent\'ee
au Bureau des Longitudes le 29 Juin, 1853.---Liouville, t.~XX. pp.~137--138. 1855.

\smallskip

\noindent\textsc{Donkin, Prof.} On a Class of Differential Equations, including those which occur in
Dynamical Problems. Part I.---Phil. Trans. 1854, pp.~71--113. Received Feb. 23rd,
read March 23rd, 1854. 1854.

\smallskip

\noindent\textsc{Donkin, Prof.} On a Class of Differential Equations, including those which occur in
Dynamical Problems. Part II.---Phil. Trans. 1855, pp.~299--358. Received Feb. 17th,
read March 22nd, 1855. 1855.

\smallskip

\noindent\textsc{Liouville}. Rapport sur un m\'emoire de M. Bouv presentant l'int\'egration des \'equations
diff\'erentielles de la M\'ecanique Analytique.---Comptes Rendus, t.~XL. p.~661, s\'eance
du 26 Mars, 1855; Liouville, t.~XX. pp.~155--160. 1855.

\smallskip

\noindent\textsc{Bouv}. Sur l'int\'egration des \'equations diff\'erentielles de la M\'ecanique Analytique
(extrait d'un m\'emoire pr\'esent\'e \`a l'Acad\'emie le 5 Mars, 1855).---Liouville, t.~XX.
pp.~165--209. 1855. [And M\'em. Savants \'Etrang. t.~XIV. pp.~33--58, 1856.]

\smallskip

\noindent\textsc{Liouville}. Note \`a l'occasion du m\'emoire pr\'ec\'edent de M. Edmund Bouv.---Liouville,
t.~XX. pp.~201--202. 1855.

\smallskip

\noindent\textsc{Brioschi}. Sopra una nuova propriet\`a degli integrali di un problema di dinamica.---
Tortolini, t.~VI. pp.~450--452. 1855.

\smallskip

\noindent\textsc{Bertrand}. M\'emoire sur quelques unes des formes les plus simples que puissent prendre
les int\'egrales des \'equations diff\'erentielles du mouvement d'un point mat\'eriel.---
Liouville, t.~II. (2$^{\rm e}$ s\'erie), pp.~113--140. 1857.

\medskip

\noindent[\textit{See} [298] \textit{at the end of papers (Vol. IV.).}]

\bigskip

\noindent\textit{2, Stone Buildings, March 28, 1856.}

\newpage


\section*{196. Note sur un Probl\`eme d'Analyse Ind\'etermin\'ee}
\addcontentsline{toc}{section}{196. Note sur un Probl\`eme d'Analyse Ind\'etermin\'ee}

\begin{center}
[From the \textit{Nouvelles Annales de Math\'ematiques} (Terquem and Gerono), t.~XVI.~(1857),
pp.~163--165.]
\end{center}

\bigskip

E\textsc{ulera} a donn\'e dans le M\'emoire: \textit{Regula facilis problemata Diophantea per numeros
integros expedite solvendi} (\textit{Commentt. Arith. Coll.}, t.~II., p.~263) la solution g\'en\'erale d'un cas
de l'\'equation ind\'etermin\'ee
\[
ax^2 + 2bx + cy + 2y = ay^2 + 2gy + d,
\]
en supposant que l'on ait la solution $x = a$, $y = b$ du nombre que
\[
aa^2 + 2ba + cy + 2 = ab^2 + 2gb + d,
\]
et en posant
\[
a^2 = ac + 1,
\]
o\`u $c$ est une quantit\'e quelconque, l'\'equation sera satisfaite par les valeurs
\begin{align*}
x &= aa + \frac{ab}{2a}(a-1)\frac{f}{a} + (a-1)g,\\
y &= -bn + a + \frac{(a-1)b}{2a} + \frac{a-1}{a}(fd).
\end{align*}
En effet, en voit sans peine que ces valeurs donnent
\[
ax^2 + 2bx + cy + 2 = a(a^2 + acb + ed) (ay^2 + 2gy + d) = 0.
\]
En rempla\c{c}ant \`a la place des coefficients $a$, $\beta$, $\gamma$, $\xi$, $\theta$ autant des nombres entiers
tels que $a_1^2$ soit un entier positif non carr\'e, on obtient une m\'ethode pour trouver
ente $r$ des nombres qui s'il existe et le coefficient sont positifs ou n\'egatifs; et si l'\'equation
$a^2 = aN + 1$ est satisfaite par $(a, a, c, q)$ et $(a', a, c, q')$, l'\'equation
\[
ax^2 + \beta x + \gamma + \xi a' + xy + \rho yx + s xy + d, \quad (1')
\]
est r\'eduite \`a la forme connue.

L'\'equation ind\'etermin\'ee (2) existe dans celle-ci
\[
(a, b, c, f, g, h)x, y_1^2 = 1, \quad \frac{(a, b, c, f, g, h)x, y, 1) = 0}{0}, \quad (2)
\]
en supposant que la forme quadratique
\[
(a, b, c, f, g, h)x, y, z)^2
\]
est l'\'equation ind\'etermin\'ee (2). En effet en supposant que la forme quadratique
\[
(a, b, c, f, g, h)x, y, z)^2 = 0
\]
se transforme en elle-m\^eme au moyen d'une substitution lin\'eaire quelconque, on peut
supposer que cette substitution soit telle que l'on ait
$x' = x = 1$ et au moment du plus $b = 0$. L'\'equation (3) se r\'eduit alors \`a montrer
que $x' = x$ aurait $de$ plus $b = 0$. L'\'equation (3) ne peut trouver par le m\'ethode pr\'ec\'edente. M. Hermite
a donn\'e cette m\'ethode dans l'ordinaire,
\[
\mathfrak{A} = ba - f\gamma, \quad \mathfrak{B} = pb - h\gamma,
\]
\[
\mathcal{B} = aa - hg^2, \quad \mathfrak{H} = ga - h\beta + h\beta.
\]

\medskip

Il faut pour cela \'ecrire
\[
ax = \mathfrak{A}^2, \quad ay = \mathfrak{B} a - p,
\]
o\`u
\[
\mathfrak{B} = \mathfrak{A}^2 + \mathfrak{B}^2 + \mathfrak{H}^2.
\]
ces \'equations sont identiquement satisfaites par
\[
(a, b, c, f, g, h)x, y, z)^2, \quad (\mathfrak{B}, b, c, f, g, h)x, y, z)^2,
\]
et un moyen de cette \'equation et de ces valeurs
\[
x' = \mathfrak{B} - \mathfrak{A}, \quad y' = \mathfrak{B} a + p,
\]
on forme tout de suite \'equation et les valeurs
\[
x = \mathfrak{B} \mathfrak{A} - p \mathfrak{B}, \quad y = \mathfrak{B} \mathfrak{A} - p \mathfrak{B} - p,
\]
en forme telles que sur les trois \'equations
\[
pa = 0, \quad \mathfrak{B} = 0 \quad \text{en ajoutant, on obtient } \mathfrak{A} = h\mathfrak{B}, \text{ce qui donne } x' = a \quad \text{et de l\`a } x' = a.
\]

Cela \'etant, les deux \'equations donnent, en rempla\c{c}ant $\mathfrak{B}$ par sa valeur,
\[
\alpha\mathfrak{A}^2 + (h - g\mathfrak{Y}) + a + g\mathfrak{B}^2 + a\mathfrak{B} = -\beta,
\]
\[
(h + g\mathfrak{Y})^2 - \alpha\mathfrak{B}^2 + a(g - g\mathfrak{X})^2 = -\beta,
\]
et en remettant pour $\mathfrak{A}$, $\mathfrak{B}$, $\mathfrak{Y}$, $\mathfrak{X}$ leurs valeurs

\noindent\textit{(page continues with formulas 4, 5, 6, 7 and conclusion)}

\medskip

\noindent\textit{London, 10 Mars, 1857.}

\newpage


\section*{197. Note on the Theory of Logarithms}
\addcontentsline{toc}{section}{197. Note on the Theory of Logarithms}

\begin{center}
[From the \textit{Philosophical Magazine}, vol.~XI.~(1856), pp.~275--280.]
\end{center}

\bigskip

An imaginary quantity $x+yi$ may always be expressed in the form
\[
x + yi = r(\cos\theta + i\sin\theta) = r e^{i\theta},
\]
where $r$ is positive, and $\theta$ is included between the limits $-\pi$ and $+\pi$. We have
in fact
\[
r = \sqrt{x^2 + y^2},
\]
and when $x$ is positive,
\[
\theta = \tan^{-1}\frac{y}{x},
\]
but when $x$ is negative,
\[
\theta = \tan^{-1}\frac{y}{x} \pm \pi,
\]
where $\tan^{-1}$ denotes an arc between the limits $-\tfrac{1}{2}\pi$, $+\tfrac{1}{2}\pi$, and where the upper
or under sign is to be employed according as $y$ is positive or negative. I use for convenience
the mark $\pm$ to denote identity of sign; we may then write
\[
\theta = \tan^{-1}\frac{y}{x} \pm \epsilon\pi,
\]
where
\[
x = +, \quad \epsilon = 0;\qquad x = -,\quad \epsilon = \pm 1 \pm y.
\]

\medskip

It should be remarked that $\theta$ has a unique value except in the single case
$x = -$, $y = 0$, where $\theta$ is indeterminately $\pm\pi$. We have, in fact, $\theta = +\pi$ or $\theta = -\pi$
according as it is considered to the limit of $x + yi$, or $x + yi$, $y$ being positive
or negative, in such an arc is the case so $r + yi$ is to be regarded.

\medskip

The preceding definition is, in fact, in the case of a positive, that given by
M. Cauchy in the \textit{Exercices de Math\'ematique}, vol.~i. [1826]; and he has there shewn that
$x$, $x'$, $xx' - yy'$ being all of them positive, the above-mentioned expression reduces itself
to zero. The general definition is that given to my \textit{M\'emoire sur quelques Formules
du Calcul Int\'egral}, Liouville, vol.~XII. [1847], p.~231 [40], but I was wrong in assuming
that the expression always reduced itself to zero. We have, in fact, in general
\[
\tan^{-1}a + \tan^{-1}\beta = \tan^{-1}\frac{a+\beta}{1-a\beta},
\]
when $1 - a\beta$ is positive; but when $1 - a\beta$ is negative (which implies that $a$, $\beta$ have
the same sign), then
\[
\tan^{-1}a + \tan^{-1}\beta = \tan^{-1}\frac{a+\beta}{1-a\beta} \pm \pi,
\]
where the upper or under sign is to be employed according as $a$ and $\beta$ are positive
or negative; or what is the same thing,
\[
\tan^{-1}a + \tan^{-1}\beta = \tan^{-1}\frac{a+\beta}{1-a\beta} \pm \epsilon\pi,
\]
where
\[
1 - a\beta = +, \quad \epsilon = 0; \qquad 1 - a\beta = -, \quad \epsilon = \pm 1 \pm a = \pm 1 \pm \beta.
\]

\medskip

This being premised, then writing
\begin{align*}
\log(x+yi) &= \log\sqrt{x^2+y^2} + i\left(\tan^{-1}\frac{y}{x} \pm \epsilon\pi\right),\\
\log(x'+y'i) &= \log\sqrt{x'^2+y'^2} + i\left(\tan^{-1}\frac{y'}{x'} \pm \epsilon'\pi\right),\\
\log(x+yi)(x'+y'i) &= \log[(xx'-yy')+i(xy'+x'y)]\\
&= \log\sqrt{(xx'-yy')^2 + (xy'+x'y)^2} + i\left(\tan^{-1}\frac{xy'+x'y}{xx'-yy'} \pm \epsilon''\pi\right),
\end{align*}
we find
\[
\log(x+yi) + \log(x'+y'i) - \log(x+yi)(x'+y'i) = (\epsilon + \epsilon' - \epsilon'')\pi i.
\]
Hence, considering the different cases:

\medskip

\noindent I.\quad $x = +,\ x' = +,\ xx' - yy' = +,$
\begin{align*}
\epsilon &= 0,\\
\epsilon' &= 0,\\
\epsilon'' &= 0,
\end{align*}
and therefore
\[
\epsilon + \epsilon' - \epsilon'' = 0.
\]

\noindent II.\quad $x = +,\ x' = +,\ xx' - yy' = -,$
\begin{align*}
\epsilon &= 0,\\
\epsilon' &= 0,\\
\epsilon'' &= \pm 1 \pm (xy' + x'y) = -\left(\frac{y}{x} + \frac{y'}{x'}\right),
\end{align*}
and therefore
\[
\epsilon + \epsilon' - \epsilon'' = \pm 1 \qquad \equiv \left(\frac{y}{x} + \frac{y'}{x'}\right).
\]

\noindent III.\quad $x = +,\ x' = -,\ xx' - yy' = +,$
\begin{align*}
\epsilon &= 0,\\
\epsilon' &= \pm 1 \pm y' = -\frac{y'}{x'} = -\frac{y}{x'},\\
\epsilon'' &= 0,\\
\epsilon'' &= \pm 1 \qquad \equiv \left(-\frac{y'}{x'}\right),
\end{align*}
and therefore
\[
\epsilon + \epsilon' - \epsilon'' = 0.
\]

\noindent IV.\quad $x = +,\ x' = -,\ xx' - yy' = -,$
\begin{align*}
\epsilon &= 0,\\
\epsilon' &= \pm 1 \pm y' = -\frac{y'}{x'},\\
\epsilon'' &= \pm 1 \pm (xy' + x'y) = -\left(\frac{y}{x} + \frac{y'}{x'}\right),
\end{align*}
and therefore
\[
\epsilon + \epsilon' - \epsilon'' = 0\qquad\text{if } \frac{y'}{x'} = -\frac{y}{x'},
\]
but
\[
\epsilon + \epsilon' - \epsilon'' = \pm 2 \qquad\text{if } \frac{y'}{x'} = -\left(\frac{y}{x} + \frac{y'}{x'}\right).
\]

\noindent V.\quad $x = -,\ x' = +,\ xx' - yy' = +,$
\begin{align*}
\epsilon &= \pm 1 \pm y = -\frac{y}{x} = -\left(\frac{y}{x}\right),\\
\epsilon' &= 0,\\
\epsilon'' &= 0,
\end{align*}
and therefore
\[
\epsilon + \epsilon' - \epsilon'' = 0.
\]

\noindent VI.\quad $x = -,\ x' = +,\ xx' - yy' = -,$
\begin{align*}
\epsilon &= \pm 1 \pm y = -\frac{y}{x},\\
\epsilon' &= 0,\\
\epsilon'' &= \pm 1 \pm (xy' + x'y) = -\left(\frac{y}{x} + \frac{y'}{x'}\right),
\end{align*}
and therefore
\[
\epsilon + \epsilon' - \epsilon'' = 0\qquad\text{if } \frac{y}{x} = -\frac{y}{x},
\]
but
\[
\epsilon + \epsilon' - \epsilon'' = \pm 2 \qquad\text{if } \frac{y}{x} = -\left(\frac{y}{x} + \frac{y'}{x'}\right).
\]

\noindent VII.\quad $x = -,\ x' = -,\ xx' - yy' = +,$
\begin{align*}
\epsilon &= \pm 1 \pm y = -\frac{y}{x},\\
\epsilon' &= \pm 1 \pm y' = -\frac{y'}{x'},\\
\epsilon'' &= 0,
\end{align*}
and therefore
\[
\epsilon + \epsilon' - \epsilon'' = 0\qquad\text{if } \frac{y}{x} = -\frac{y'}{x'},
\]
but
\[
\epsilon + \epsilon' - \epsilon'' = \pm 2 \qquad\text{if } \frac{y}{x} = \frac{y'}{x'}.
\]

\noindent VIII.\quad $x = -,\ x' = -,\ xx' - yy' = -,$
\begin{align*}
\epsilon &= \pm 1 \pm y = -\frac{y}{x},\\
\epsilon' &= \pm 1 \pm y' = -\frac{y'}{x'},\\
\epsilon'' &= \pm 1 \pm (xy' + x'y) = -\left(\frac{y}{x} + \frac{y'}{x'}\right),
\end{align*}
and therefore
\[
\epsilon + \epsilon' - \epsilon'' = \pm 1 \qquad \equiv -\left(\frac{y}{x} + \frac{y'}{x'}\right).
\]

\medskip

Hence writing
\[
\log(x+yi) + \log(x'+y'i) - \log(x+yi)(x'+y'i) = E\pi i,
\]
we have $E = 0$, except in the following cases, viz.

\medskip

\noindent 1.\quad (See IV.) \quad $x = +,\ x' = -,\ xx' - yy' = -, \quad \frac{y'}{x'} = -\left(\frac{y}{x} + \frac{y'}{x'}\right),$
\quad where
\[
E = \pm 2 \pm \left(\frac{y}{x} + \frac{y'}{x'}\right).
\]

\noindent 2.\quad (See VI.) \quad $x = -,\ x' = +,\ xx' - yy' = -, \quad \frac{y}{x} = -\left(\frac{y}{x} + \frac{y'}{x'}\right),$
\quad where
\[
E = \pm 2 \pm \left(\frac{y}{x} + \frac{y'}{x'}\right).
\]

\noindent 3.\quad (See VII.) \quad $x = -,\ x' = -,\ xx' - yy' = +, \quad \frac{y}{x} = \frac{y'}{x'},$
\quad where
\[
E = \pm 2 \pm \frac{y}{x}.
\]

\noindent 4.\quad (See VIII.) \quad $x = -,\ x' = -,\ xx' - yy' = -,$
\quad where
\[
E = \pm 2 \mp \left(\frac{y}{x} + \frac{y'}{x'}\right).
\]

\medskip

It thus appears that when any one of the arguments $x+yi$, $x'+y'i$ are real and negative, or
in the imaginary cases when the product of the proper expression for the
arguments must be positive but their sums must be negative, then the equation
\[
\log(x+yi) + \log(x'+y'i) = \log(x+yi)(x'+y'i)
\]
ceases to hold in any case.

\medskip

The preceding results do not apply to the case where any one of the arguments
$x+yi$, $x'+y'i$, $(x+yi)(x'+y'i)$ is real and negative, in which case the formula
$\log z = +\pi i$ being still negative, has in addition to its arbitrary positive
or negative value $\pm\pi i$, what is the same thing the logarithm is absolutely indeterminate.
The general methods are illustrated by applications to the problem of
three points of rotation; the former problem is specially discussed
in section 7., but the results obtained (and which, to the author remarks, afford an
example of the so-called ``elimination of the nodes'') do not come within the plan of
the present report.

\bigskip

\noindent\textit{2, Stone Buildings, March 18, 1856.}

\newpage


\section*{198. Note upon a Result of Elimination}
\addcontentsline{toc}{section}{198. Note upon a Result of Elimination}

\begin{center}
[From the \textit{Philosophical Magazine}, vol.~XI.~(1856), pp.~378--379.]
\end{center}

\bigskip

I\textsc{f} the quadratic function
\[
(a,\ b,\ c,\ f,\ g,\ h)(x,\ y,\ z)^2
\]
break up into factors, then representing one of these factors by $\xi x + \eta y + \zeta z$, and taking
any arbitrary quantities $a$, $\beta$, $\gamma$, the factor in question, and therefore the quadratic
function is reduced to zero by substituting $\beta\zeta - a\eta$, $\gamma\xi - a\zeta$, $a\eta - \beta\xi$ in the place of
the ideas, a more general; the function then results of the operation of performed on
$U$, and $QPU$ denotes the result of the operation $Q$ performed on $PU$, and generally
in such combinations of symbols, each operation is considered as affecting the operand
denoted by means of all the symbols on the right of the operation in question. Now
considering the operation $QPU$ it is easy to see that we may write
\[
QPU = (Q + P)U + (QP)U,
\]
where on the right-hand side $(Q + P)$ and $(QP)$ signify as follows: viz. $Q + P$ denotes
the new algebraical product of $Q$ and $P$, while $QP$ (consistently with the general
notation as before explained) denotes the result of the operation $Q$ performed upon
$P$ as operand; and the two parts $(Q + P)U$ and $(QP)U$ denote respectively the
results of the operations $(Q + P)$ and $(QP)$ performed each of them upon $U$ as operand.
It is proper to remark that $(Q + P)$ and $(P+Q)$ have precisely the same meaning,
and the symbol may be written in either form indifferently. But without a more
convenient notation, it would be difficult to find the corresponding expression for
$RQPU$, \&c. This, however, has at one effected by means of the analytical forms
called \textit{trees} (see figs. 1, 2, 3), which contain all the trees which can be formed with
one branch, two branches, and three branches respectively.

\medskip

\noindent[Figures of trees with single branches and combinations are described in the text.]

\medskip

The coefficients $a$, $b$, $c$, $d$ are linear functions of coordinates, the equation
\[
aa-b^2 = 0,\quad bd - cd = 0,\quad bd - c^2 = 0
\]
(equivalent, of course, to two equations) belong to a curve of the third order in space;
the edge of regression of the developable surface obtained by putting the discriminant
equal to zero, or which has for its apolar
\[
-x^3d^2 + 6abcd - 4ac^3 - 4b^3d + 3b^2c^2 = 0,
\]
and, moreover, the above forms are the general representations of any curve of the
third order, and developable surface of the fourth order. The question arises, to find
the equation of the cone of the third order having an arbitrary point for its vertex,
and passing through the curve of the third order. This may be done by Joachimsthal's
method (vide. p.~21, q., f. be the values of $a$, $b$, $c$, $d$ at the vertex of the cone, and is
the equations of the cone, say in
\[
aa - b^2 = 0,\quad bd - c^2 = 0,
\]
for $a$, $b$, $c$, $d$ write $aa + ka$, $ab + kb$, $ac + kc$, $ad + kd$; if then the equations so obtained
in $k$ are eliminated, the resulting equation will be that of the cone of the third order.

\bigskip

\noindent\textit{2, Stone Buildings, May 1, 1856.}

\newpage


\section*{199. Note on the Theory of Elliptic Motion}
\addcontentsline{toc}{section}{199. Note on the Theory of Elliptic Motion}

\begin{center}
[From the \textit{Philosophical Magazine}, vol.~XI.~(1856), pp.~425--428.]
\end{center}

\bigskip

I\textsc{f}, as usual, $r$, $\dot{r}$ denote the radius vector and longitude, and $\mu$ the central mass,
then the Vis Viva and Force function are respectively
\[
T = \tfrac{1}{2}(\dot{r}^2 + r^2\dot{\theta}^2),
\]
\[
U = \frac{\mu}{r},
\]
and writing
\[
\frac{dT}{d\dot{r}} = p,\qquad r^2\dot{\theta} = p,
\]
\[
\frac{dT}{d\dot{\theta}} = r^2\dot{\theta} = q,
\]
we have $\dot{r}^2 = p$, $\dot{\theta}^2 = q$, and $T = \tfrac{1}{2}\left(p^2 + \dfrac{q^2}{r^2}\right)$, whence, putting $H = T - U$, the value of
$H$ is
\[
H = \tfrac{1}{2}\left(p^2 + \frac{q^2}{r^2}\right) - \frac{\mu}{r};
\]
and by Sir W. R. Hamilton's theory, the equations of motion are
\[
\frac{dr}{dt} = \frac{dH}{dp},\quad \frac{d\theta}{dt} = \frac{dH}{dq},
\]
\[
\frac{dp}{dt} = -\frac{dH}{dr},\quad \frac{dq}{dt} = -\frac{dH}{d\theta},
\]
or substituting for $H$ its value, the equations of motion are
\[
\frac{dr}{dt} = p,
\]
\[
\frac{d\theta}{dt} = \frac{q}{r^2},
\]
\[
\frac{dp}{dt} = \frac{q^2}{r^3} - \frac{\mu}{r^2},
\]
\[
\frac{dq}{dt} = 0.
\]

Putting, as usual, $p = a^2(1 - e^2)$, and introducing the eccentric anomaly $u$, which is given
as a function of $t$ by means of the equation
\[
nt + \xi = u - e\sin u,
\]
$\left(\text{so that }\dfrac{du}{dt} = \dfrac{n}{1 - e\cos u}\right)$, the integral equations are
\[
q = na^2\sqrt{1 - e^2},
\]
\[
p = \frac{ena\sin u}{1 - e\cos u},
\]
\[
r = a(1 - e\cos u),
\]
\[
\theta = \varpi + \tan^{-1}\left(\frac{\sqrt{1-e^2}\sin u}{\cos u - e}\right),
\]
where the constants of integration $a$, $e$, $\varpi$ denote as usual the mean distance, the
eccentricity, the mean anomaly at epoch, and the longitude of perisaster.

Suppose that $q_0$, $p_0$, $r_0$, $\theta_0$, $u_0$ correspond to the time $t_0$ ($q$ is constant, so that
$q_0 = q$), and write
\[
V = nt'(u - u_0 + e\sin u - e\sin u_0),
\]
joining to this the equations
\[
r = a(1 - e\cos u),\quad r_0 = a(1 - e\cos u_0),
\]
\[
\theta - \varpi = \tan^{-1}\left(\frac{\sqrt{1-e^2}\sin u}{\cos u - e}\right) - \tan^{-1}\left(\frac{\sqrt{1-e^2}\sin u_0}{\cos u_0 - e}\right),
\]
$u$, $u_0$, $\theta$ will be functions of $a$, $e$, $r$, $r_0$, $\theta$, and consequently $\xi$ being throughout considered
as a function of $a$, $e$, $\xi$, $H$, in this view, the characteristic function $V$ of
Sir W. R. Hamilton, and according to his theory we ought to have
\[
dV = \tfrac{1}{2}n(u - u_0)da + p\,dr - p_0\,dr_0 + q\,d\theta - q_0\,d\theta_0.
\]

To verify this, I form the equation
\begin{align*}
dV &= [nu(u - u_0 + e\sin u - e\sin u_0)\,da\\
&\quad + nt'\{(1 - e\cos u)\,du - (1 - e\cos u_0)\,du_0 - \sin u\,de + \sin u_0\,de\}]\\
&= \frac{n\sin u}{1 - e\cos u}\left[de - (1 - e\cos u)\,da - r\sin u\,du + a\sin u_0\,de\right]\\
&\qquad + na\sqrt{1 - e^2}\left[d\theta - \frac{1}{\sqrt{1-e^2}(1 - e\cos u)}((1 - e^2)\,du + ae\sin u\,de)\right.\\
&\qquad - dH + \left.\frac{1}{\sqrt{1-e^2}(1 - e\cos u_0)}((1 - e^2)\,du_0 + ae\sin u_0\,de)\right],
\end{align*}
the coefficient of $da$ on the right-hand side is
\[
nu'(1 + e\cos u) - \frac{nut'\sin^2 u}{1 - e\cos u} - \frac{ane^2(1 - e^2)}{1 - e\cos u},
\]
which vanishes, and similarly the coefficient of $du$, also vanishes; the coefficient of $de$
is the difference of two parts, the first of which is
\[
\frac{ne\sin u}{1 - e\cos u}\,de - na\sin u\,e^2 du,
\]
and the second is the form of the form $u$, by virtue of which the coefficient therefore is
\[
\frac{1}{2}n(u - u_0 + e\sin u - e\sin u_0).
\]
We have therefore
\[
dV = \tfrac{1}{2}n(u - u_0 + e\sin u - e\sin u_0)\,da
\]
\[
+ \frac{ne\sin u}{1 - e\cos u}\,dr - na\sqrt{1 - e^2}\,d\theta,
\]
\[
- \frac{ne\sin u_0}{1 - e\cos u_0}\,dr_0 - na\sqrt{1 - e^2}\,d\theta_0,
\]
or what is the same thing,
\[
dV = \tfrac{1}{2}n(u - u_0)\,da + p\,dr - p\,d\theta\,dr_0 - q\,d\theta_0 - q\,d\theta_0,
\]
the equation which was to be verified.

\bigskip

\noindent\textit{2, Stone Buildings, March 26, 1856.}

\newpage


\section*{200. On the Cones which pass through a given curve of the third order in space}
\addcontentsline{toc}{section}{200. On the Cones which pass through a given curve of the third order in space}

\begin{center}
[From the \textit{Philosophical Magazine}, vol.~XII.~(1856), pp.~20--22.]
\end{center}

\bigskip

T\textsc{he} following investigation is connected with the theory of the cubic
\[
(a, b, c, d)(x, y)^3,
\]
and in particular with a theorem that the determinant formed with the second
differential coefficients of the discriminant gives the square of the discriminant.

Consider the coefficients $a$, $b$, $c$, $d$ as linear functions of coordinates, the equations
\[
ac - b^2 = 0,\quad bd - cd = 0,\quad bd - c^2 = 0
\]
(equivalent, of course, to two equations) belong to a curve of the third order in space;
the edge of regression of the developable surface obtained by putting the discriminant
equal to zero, or which has for its apolar
\[
-a^2d^2 + 6abcd - 4ac^3 - 4b^3d + 3b^2c^2 = 0,
\]
and moreover, the above forms are the general representations of any curve of the
third order, and developable surface of the fourth order. The question arises, to find
the equation of the cone of the third order having an arbitrary point for its vertex,
and passing through the curve of the third order. This may be done by Joachimsthal's
method: let $a$, $b$, $c$, $d$ be the values of $a$, $b$, $c$, $d$ at the vertex of the cone, and in
the equations of the cone, viz.
\[
ac - b^2 = 0,\quad ad - bc = 0,\quad bd - c^2 = 0,
\]
for $a$, $b$, $c$, $d$ write $a + ka$, $b + kb$, $c + kc$, $d + kd$; if then the equations so obtained
in $k$ are eliminated, the resulting equation will be that of the cone of the third order.

The substitution in question gives
\[
La + Ma = R + S' a + 0,
\]
\[
La' + 2Ma = M' a + 0,
\]
where
\begin{align*}
L &= \tfrac{1}{2}(ac - b^2),& M &= -ac + bd - 2hd,& N &= \tfrac{1}{2}(by - 8'),\\
L &= \tfrac{1}{2}(ad - c),& M &= bd - bd - cd,& N &= \tfrac{1}{2}(ed - dc),\\
L'' &= \tfrac{1}{2}3M''M + 2g',& M'' &= bd - bd + cd - 3M'',& X'' &= \tfrac{1}{2}(cd - bb);
\end{align*}
the result of the elimination is
\[
4(LX' - M)(LX'' - 2M'') - (LX'' + LX' - 2MM'')^2 = 0,
\]
and we have
\[
LX - M^2 = 4(a-b)(c-bd) - (4-bc)^2 = ac + 3b'(a'-b'),
\]
\[
LX'' - M'' = 4(a-d)(bd - ce) - (ad - bc)^2 = (ad - bc)^2,
\]
\[
LX'' + LX' - 2MM'' = 8(ad - bc)(bd - cd) + (4 - bd)(ad - bc) + 8M.M,
\]

Write for shortness
\[
by - cf = i, \quad cn - dy = m, \quad cd - bn = n, \quad ad - cf = w,
\]
\[
ad - by = w, \quad bd - cf = u, \quad sf - by = w, \quad dd - 2f = y,
\]
values which (if $a$, $b$, $c$, $d$ are linear functions of coordinates $x$, $y$, $z$, $w$) might, $u$, $v$, $w$
this is equal to
\[
-16(P - cf - syh + day - 4)bw + uw^2,
\]
or, substituting and reducing,
\[
P - cy h + lay = 1bw - 1ux + b'y w - 3a'(P_x - dy x) + 3y(d'a - 4'8) bd,
\]
\[
+ k(P - 4y k) + 3by^2 + by^2 + bd\\
- (a8) k\Omega_b'^2 + 8y'd + 4cd + 8cd,
\]

I do not however refer the result obtained.

\medskip

Now putting
\[
bd - c^2 = p,\qquad bc - cd = q,\qquad ac - b^2 = r,
\]
and in like manner
\[
8d - q' = D,\qquad By - cb = Q,\qquad ay - bg = R,
\]
and reducing, the final result may be expressed in the form
\begin{align*}
P[\quad &= 2Dpq + (pn - 2)q + p(2a - py)r]\\
+ Q[\quad &-Wn + ((ph - 2)r - py)\quad He]\\
+ R[\quad & + (Dy - 2(b)p - (Db - Ny) + 0]\\
&= 0,
\end{align*}
where $a$, $b$, $c$, $d$ are current coordinates, and $p$, $q$, $r$ are quadratic functions of $a$, $b$, $c$, $d$,
the equation $P = 0$ of the second order in which the vertex (current point) of the cone $(p$, $q$, $r$ ) cone is the quadratic functions of $a$, $b$, $c$, $d$,
the equation of the cone is in $p$, $q$, $r$, $R$, $r$ are quadratic; and note VII
this is equal to (suppose) is
\[
F = a + ay + a a v = R,
\]
of the cone in question (so an arbitrary quantity), is
\[
p + uq + sP'r = 0,
\]
or so all length quantity
\[
(bd - c^2) + (a + xc)(c - cd) + 8c^2(ac - b^2) = 0,
\]
in fact this equation is evidently that of a surface of the second order passing
through the curve; and there is no difficulty in shewing that it is a cone.

\bigskip

\noindent\textit{2, Stone Buildings, May 1, 1856.}

\newpage


\section*{201. Second Note on the Theory of Logarithms}
\addcontentsline{toc}{section}{201. Second Note on the Theory of Logarithms}

\begin{center}
[From the \textit{Philosophical Magazine}, vol.~XII.~(1856), pp.~354--363.]
\end{center}

\bigskip

T\textsc{he} theory of logarithms, as developed in my first note (\textit{Phil. Mag.} April 1856), [197]
may be exhibited in a clearer light by considering, instead of $\log a + \log b - \log ab = E\pi i$,
the equation $\log \dfrac{x}{a} = \log b + \log a - E\pi i$, in which more readily enables the answering
\textit{\`a priori} for the discontinuity in the values of $E$. Writing then for $b$ the complex
values $x + yi$, $c + yi$, we have
\[
\log\frac{x'+y'i}{x+yi} = \log(x'+y'i) - \log(x+yi) + E\pi i,
\]
where, according to the assumed definition of a logarithm,
\[
\log(x+yi) = \log\sqrt{x^2+y^2} + i\left(\tan^{-1}\frac{y}{x} \pm \epsilon\pi\right),
\]
in which $\tan^{-1}$ denote an arc between the limits $-\tfrac{1}{2}\pi$, $+\tfrac{1}{2}\pi$, and $\tan^{-1}\dfrac{y}{x}$ is an arc between
the limits $-\tfrac{1}{2}\pi$, $+\tfrac{1}{2}\pi$. The coefficient $\epsilon$ is equal to zero when $x$ is positive; but
when $x$ is negative, then $\epsilon = +1$ or $-1$, according as $y$ is positive or negative, i.e. we
have
\begin{align*}
x &= +, & \epsilon &= 0,\\
x &= -, & \epsilon &= \pm 1\pm y;
\end{align*}
and of course the other logarithm in the equation have an analogous signification.

Hence, attending to the equation
\[
\tan^{-1}a + \tan^{-1}\alpha = \tan^{-1}\frac{\beta+a}{1+a\beta}\pm\epsilon''\pi,
\]
where, when $1 + a\beta$ is positive, $\epsilon''$ is equal to zero; but when $1 + a\beta$ is negative,
$\epsilon''$ is equal to $+1$ or $-1$, according as $\beta + a$ is positive or negative; or what is the
same thing, ($\beta$ being of opposite signs when $1 + a\beta$ is negative), or to
\begin{align*}
1 + a\beta &= +, & \epsilon'' &= 0;\\
1 + a\beta &= -, & \epsilon'' &= \pm 1 \pm a = \mp 1 \mp \beta.
\end{align*}
we find
\[
E = \epsilon + \epsilon' - \epsilon'',
\]
where $\epsilon$, $\epsilon'$, $\epsilon''$ are defined by the conditions
\begin{align*}
x &= +, & \epsilon &= 0,\\
x &= -, & \epsilon &= \pm 1 \pm y,\\
x' &= +, & \epsilon' &= 0,\\
x' &= -, & \epsilon' &= \pm 1 \pm y',\\
xx' + yy' &= +, & \epsilon'' &= 0,\\
xx' + yy' &= -, & \epsilon'' &= \pm 1 \pm xy' = \mp 1 \mp x'y.
\end{align*}

Suppose, to fix the ideas, that $x$, $y$ are each of them positive; we have $x = 0$,
and considering the several cases:

\medskip

\noindent 1.\quad $x' = +,\ y' = +,$

Here $xx' + yy' = +$, $1 + \frac{y}{x}\frac{y'}{x'} = +$; and consequently not only $\epsilon' = 0$, but also $\epsilon'' = 0$,
$\epsilon'' = 0$, and thence $E = 0$.

\noindent 2.\quad $x' = +,\ y' = -,$

Here $\epsilon = 0$. Moreover, $xx'$ being positive, and $\epsilon yy'$ and $1 + \frac{y}{x}\frac{y'}{x'}$ will have the
same sign. If they are both positive, $xx' + yy'$ is positive, and $1 + \frac{y}{x}\frac{y'}{x'}$ is negative.
then $x' = 1 \pm xy' = +(\pm \frac{y}{x} - \frac{y'}{x'})$ (since $xy'$ is negative) = $-\frac{y}{x}$, i.e. $\epsilon'' = +1$, and
$1 + \frac{y'}{x'}\frac{y'}{x'}$ being positive, we have $\epsilon'' = 0$. Hence in each case $E = 0$.

\noindent 3.\quad $x' = -,\ y' = +,$

Here $\epsilon' = \pm 1 \pm y'$, i.e. $\epsilon = -1$. Also $xx'$ and $yy'$ being each negative, $xx' + yy'$
will be negative, and therefore
\[
\epsilon'' = \pm 1 \pm xy' = -\left(\frac{y'}{x'}\right),
\]
i.e. if $\frac{y'}{x'} > \frac{y}{x}$, then $\epsilon'' = -1$; but if $\frac{y'}{x'} < \frac{y}{x}$, then $\epsilon'' = +1$. $1 + \frac{y}{x}\frac{y'}{x'}$ being positive,
$\epsilon'' = 0$. Hence if $\frac{y'}{x'} > \frac{y}{x}$, $E - 1 - 1 = 0$; but if $\frac{y'}{x'} < \frac{y}{x}$, then $E = 1 + 1 = 2$.

\noindent 4.\quad $x' = -,\ y' = -,$

Here $x' = \pm 1 \pm y'$, i.e. $\epsilon' = -1$. Also $xx'$ and $yy'$ being each negative, $xx' + yy'$
will be negative, and therefore
\[
\epsilon'' = \pm 1 \pm xy' = -\left(\frac{y'}{x'}\right),
\]
i.e. $\epsilon'' = -1$. Hence $1 + \frac{y}{x}\frac{y'}{x'} = -$. If $\frac{y'}{x'} > \frac{y}{x}$, $\epsilon'' = 0$. Hence in such case $E = 0$.

\smallskip

Consider $(x, y)$ $(x', y')$ as the rectangular coordinates of two points $A$, $A'$. In
the case which has been considered, the point $A$ has been taken in the positive
quadrant; and the preceding discussion shews that we have always $E = 0$, except in
the case in which the negative portion of the axis $x$. But when this happens, then if
the line $AA'$, considered as drawn from $A$ to $A'$, passes from above to below the
axis of $x$, we have $E = +2$; but if the line $AA'$, considered as drawn from $A$ to $A'$,
passes from below to above the axis of $x$, then $E = -2$. So that treating the points
$A$, $A'$ as the geometrical representations of two values $x + yi$, $x' + y'i$, we have for the
discontinuity in the value of the difference of the logarithms in passing in a continuous
quadratic; and the preceding discussion shews that we have always $E = 0$, except in
the case where the finite line $AA'$ meets the negative portion of the axis of $x$, in
which case we have $E = +2$ or $E = -2$. The same thing is true generally in whichever quadrant

\smallskip

$A$ is situated, i.e. we have always $E = 0$, except in the case in which the finite line
$AA'$ meets the negative portion of the axis of $x$. But when this happens, then if
the line $AA'$, considered as drawn from $A$ to $A'$, passes from above to below the
axis of $x$, we have $E = +2$; but if the line $AA'$, considered as drawn from $A$ to $A'$,
passes from below to above the axis of $x$, then $E = -2$.

\medskip

It thus appears that when the values $x+yi$, $x'+y'i$ are exhibited by means of two
points $A$, $A'$ in a plane, the discontinuity in the value of $E$ depends on the path of
the line which joins the two points. So that treating the points $A$, $A'$ as the geometrical
representations of the values $x+yi$, $x'+y'i$, the difference of the logarithms is
the result of integration $\int_A^{A'}\dfrac{dz}{z}$, the path of integration being the line from $A$ to $A'$.
Consider in general the definite integral
\[
\int_A^{A'} d u \, \phi u,
\]
where $\phi$, $\phi'$ are complex numbers of the form $x+yi$, $x'+y'i$; and take $A$, $A'$ the geometrical
representations of these limits, and the variable point $P$ as the geometrical
representation of the complex variable $z$. The value of the definite integral will
depend in general not only upon the values which $\phi$, $\phi'$ assume as constants, but also
upon the path which the variable $z$ supposes to be successively assumed in passing from
$A$ to $A'$.

\medskip

To order therefore to give a precise signification to the notation, we must for the
path of the point $P$, and it is natural to assume that the path is a right line (of
course there exist also an infinity of paths which give the same value to the definite
integral, or as we may call them, paths equivalent to the right line); but the consideration
of these would be a tedious complication of the definition, and it is better to attend to the
single path---the right line). The definition in an exact sense
executed into an analytical mode of computing the negative part of the axis of $x$ in
which case we have $E = +2$ or $E = -2$ (according as the line goes from above to
below or below to above the axis), then there may be considered as a more general
spherical mass having a double natural connection of integration without limits. (For
which is in fact identical with that between $\phi$ and $\phi'$, and in fact the equation between
$\theta$, $\theta'$ is
\[
\tan\theta\tan\phi = \frac{1}{R},
\]
which, if $\phi = \cos a$, gives $\theta = \sin (R - r)$.

\smallskip

Richelot has shewn, by precisely similar reasoning, that for circles of the sphere
we have
\[
\frac{d\phi}{\sqrt{\cos^2 r - (\cos R \sin u + \sin R \cos u \cos 2\phi)^2}} = \frac{d\phi'}{\sqrt{\cos^2 r - (\cos R \sin u + \sin R \cos u \cos 2\phi')^2}},
\]
which is of the form
\[
\frac{d\phi}{\sqrt{1 - (k + \mu \cos 2\phi)^2}} = \frac{d\phi'}{\sqrt{1 - (k + \mu \cos 2\phi')^2}},
\]
where
\[
k = \frac{\cos R \sin u}{\cos r}, \qquad \mu = \frac{\sin R \cos u}{\cos r}.
\]

\bigskip

\noindent\textit{2, Stone Buildings, September 13, 1856.}

\newpage


\section*{202. Supplementary Remarks on the Porism of the In-and-Circumscribed Triangle}
\addcontentsline{toc}{section}{202. Supplementary Remarks on the Porism of the In-and-Circumscribed Triangle}

\begin{center}
[From the \textit{Philosophical Magazine}, vol.~XIII.~(1857), pp.~19--30.]
\end{center}

\bigskip

I\textsc{n} my former papers (see \textit{Phil. Mag.} August and November, 1853, [115, 116]) I
established (as part of a more general one) the following theorem, viz. the condition
that there may be inscribed in the conic $U = 0$ an infinity of triangles circumscribed
about the conic $V = 0$, is that if we develope in ascending powers of $k$ the square root
of the discriminant of $kU + V$, the coefficient of $k$ in this development must vanish.
Thus writing
\[
\operatorname{disc.}(kV + U) = (K, \mathfrak{R}, \mathfrak{R}', K')(k, 1)^3,
\]
the condition in question is found to be
\[
36\mathfrak{R}^2 - 4KW = 0.
\]

The following investigations, although relating only to particular cases of the
theorem, are, I think, not without interest. If the equation of the conic containing
the angles is
\[
U = \xi a y x + \mathfrak{B} c x + \mathfrak{C} a y = 0,
\]
and the equation of the conic touched by the sides is
\[
V = a^2 x^2 + b^2 y^2 + c^2 z^2 - 2bc x z - 2ca y x - 2ab x y = 0,
\]
that is,
\[
K = -A,\qquad \mathfrak{R} = -\tfrac{1}{3}(a + b + c),\qquad \mathfrak{R}' = -\tfrac{1}{3}(a + b + c)^2,\qquad K' = \text{disc.}.
\]

\medskip

\noindent[\textit{The detailed investigation extends through the analysis of porism
conditions for the in-and-circumscribed triangle, involving determinants, conics,
and the formulae of Richelot for the sphere.}]

\medskip

By comparing with the corresponding formula \textit{in plano}, we arrive at Richelot's
conclusion, that the formul\ae{} for the sphere may be deduced from those \textit{in plano} by
writing in the place of $\dfrac{R}{r}$, $\dfrac{a}{r}$, the functions $\dfrac{\tan R}{\tan r}$, $\dfrac{\sin a}{\cos R \sin r}$, respectively.

\bigskip

\noindent\textit{2, Stone Buildings, October 1, 1856.}

\newpage


\section*{203. On the Theory of the Analytical Forms called Trees}
\addcontentsline{toc}{section}{203. On the Theory of the Analytical Forms called Trees}

\begin{center}
[From the \textit{Philosophical Magazine}, vol.~XIII.~(1857), pp.~172--176.]
\end{center}

\bigskip

A \textsc{symbol} such as $AU_0 + BU_1 + \ldots$, where $A$, $B$, \&c. contain the variables $x$, $y$, \&c. in
respect to which the differentiations are to be performed, partakes of the natures
of an operand and operator, and may be therefore called an Operandator. Let
$P$, $Q$, $R$, \&c. be any operandators, and let $U$ be a symbol of the same kind, or to
the ideas a mere operand; $PU$ denotes the result of the operation of $P$ performed on
$U$, and $QPU$ denotes the result of the operation $Q$ performed on $PU$; and generally
in such combinations of symbols, each operation is considered as affecting the operand
denoted by means of all the symbols on the right of the operation in question. Now
considering the operation $QPU$, it is easy to see that we may write
\[
QPU = (Q + P)U + (QP)U,
\]
where on the right-hand side $(Q + P)$ and $(QP)$ signify as follows: viz. $Q + P$ denotes
the mere algebraical product of $Q$ and $P$, while $QP$ (consistently with the general
notation as before explained) denotes the result of the operation $Q$ performed upon
$P$ as operand; and the two parts $(Q + P)U$ and $(QP)U$ denote respectively the
results of the operations $(Q + P)$ and $(QP)$ performed each of them upon $U$ as operand.
It is proper to remark that $(Q + P)$ and $(P + Q)$ have precisely the same meaning,
and the symbol may be written in either form indifferently. But without a more
convenient notation, it would be difficult to find the corresponding expression for
$RQPU$, \&c. This, however, has been effected by means of the analytical forms
called \textit{trees}.

\medskip

The inspection of these figures will at once shew what is meant by the term in
question, and by the terms root, branches (which may be either main branches, intermediate,
or free branches), and knots (which may be either the root itself,
or proper knots, or the extremities of the free branches). To apply this to the
question in hand, $PU$ consists of a single term, represented by fig. 1 ($bis$); $QPU$
consists, as we have seen, of two terms represented by the two parts of fig. 2 ($bis$); the
first part represents the term $(Q + P)U$, and the second part represents the term
$(QP)U$. And it is obvious then fig. 2 ($bis$) is at once formed from the figure 1
($bis$) by adding on a branch terminated by $Q$ at each of the knots: the upper part
of fig. 1 ($bis$). In like manner $RQPU$ consists of six terms represented by the six
parts of fig. 3 ($bis$), and this figure is at once formed from fig. 2 ($bis$) by adding
on a branch terminated by $R$ at each of the knots of each part of fig. 2 ($bis$). It is
hardly necessary to remark that the first part of fig. 3 ($bis$) denotes what, in the
notation first explained, would be denoted by $(R + Q + P)U$, the second term what
would be like manner be denoted by $(RQ + P)U$, and so on; the last part being the
term which would be denoted by $((RQ)P)U$, viz. the operation upon $Q$, giving the
operandator $RQ$, which operates upon $P$, giving the operandator $(RQ)P$, which finally
operates upon $U$.

\medskip

The figures 1 ($bis$), 2 ($bis$), 3 ($bis$), contain the same trees as are contained in the
corresponding figures 1, 2, 3, \&c.; only, on account of the different modes of filling up,
trees are considered as so many distinct trees in a figure of the second set which
correspond to one and the same tree in the corresponding figure of the first set. A
difference in the number of trees first occurs in the figures 3 and 3 ($bis$), the
number of trees in figure 3 being 4, while the number of trees in figure 3 ($bis$) is 9.
Again, on account of similarity, only the figures of the first set are considered as truly
distinct from each other; thus, in fig. 3 the first tree is fill up in two way, viz.
the third tree of fig. 3 ($bis$); in like manner $U$ at the root and $R$, $Q$, $P$ at the other
knots in any one many ways; but the trees so obtained are considered as one and the same
tree, for, then we must only fill up the first tree of fig. 3 ($bis$) in the order of fig. 3
(which figure for shortness in considered as one of fig. 3 ($bis$)). Again, in fig. 3, the
second tree contains the latter contains six trees, viz. the first tree,
the second, third and fourth tree, the fifth tree and the sixth tree of fig. 3 ($bis$)
correspond respectively to the first tree, the second tree, the third tree, and the fourth
tree of fig. 3, \&c. To derive fig. 3 ($bis$) from fig. 3, we must fill up the trees of fig.
3 with $U$, $R$, $Q$, $P$ in all the ways possible: thus the four orders of fig. 3 of
with $U$, $R$, $Q$, $P$, may be filled up in one way; the second of fig. 3 ($bis$)
correspond respectively to the first tree, the second tree, the third tree, and the fourth
tree of fig. 3, \&c. To derive fig. 3 ($bis$) from fig. 3, we must fill up the trees of fig.
3 with $U$, $R$, $Q$, $P$ in all the ways possible: the only restriction is that the root
to be filled in every way, but the trees so obtained are considered as one and the same
tree, and there may only be no restriction in either form indifferently. The number of trees in
the resulting fig. 3 ($bis$) which may then be filled in any direct manner are those in
the annexed figures, and the number is therefore two. It is not difficult to
see that we have in this case ($R_r$ being the number of such trees with $r$ free branches),
\[
(1 - x)^{-1}(1 - x^2)^{-R_1}(1 - x^3)^{-R_2}\cdots = 1 + x + 2R_2 x^2 + 2R_3 x^3 + 2R_4 x^4 + \&\textrm{c.};
\]

\medskip

To obtain the law for the number $A_n$ of the trees with $n$ branches, we consider how
the trees with $n$ branches can be formed by means of those of a smaller number of
branches. A tree with $n$ branches has either a single main branch, or else two main
branches, three main branches, etc.\ to $n$ main branches. If the tree has one main
branch, it can only be formed by adding on to this main branch a tree with $(n-1)$
branches, i.e. $A_n$ contains a part $A_{n-1}$. If the tree has two main branches, then
$p + q$ being a partition of $n - 2$, the tree can be formed by adding on to one main
branch a tree of $p$ branches, and to the other main branch a tree of $q$ branches; the
number of trees so obtained is $A_p A_q$; this, however, assumes that the parts $p$ and $q$
are unequal; if they are equal, it is easy to see that the number of trees is only
$\tfrac{1}{2}A_p(A_p + 1)$. Hence $p + q$ being any partition of $n - 2$, $A_n$ contains the part $A_p A_q$,
if $p$ and $q$ are unequal, and the part $\tfrac{1}{2}A_p(A_p + 1)$ if $p$ and $q$ are equal. In like
manner, considering the trees with three main branches, then if $p + q + r$ is any
partition of $n - 3$, $A_n$ contains the part $A_p A_q A_r$ if $p$, $q$, $r$ are unequal; the part
$\tfrac{1}{2}A_p(A_p + 1)A_r$ if $p$ and $q$ are equal, and so on, until lastly we have a single
tree with $n$ main branches, or $A_n$ contains the part $A_1$. A little consideration will
shew that the preceding rule for the formation of the number $A_n$ is completely expressed
by the equation
\[
(1 - x)^{-A_1}(1 - x^2)^{-A_2}(1 - x^3)^{-A_3}\cdots = 1 + A_1 x + A_2 x^2 + A_3 x^3 + A_4 x^4 + \&\textrm{c.},
\]
and consequently that we may, by means of this equation, calculate successively for
the different values of the number $A_n$ of the trees with $n$ branches. The calculation
may be effected very easily as follows: [the table as originally printed contained at the
end of it some errors of calculation which were corrected, B.A. Report for 1875, p.~258].
\[
\begin{array}{c|*{11}{c}|c}
A_1 = & (1) & 1 & 1 & 1 & 1 & 1 & 1 & 1 & 1 & 1 & 1 & (1-x)^{-1}\\
A_2 = & 1 & (1) & 2 & 2 & 3 & 3 & 4 & 4 & 5 & 5 & 6 & (1-x^2)^{-1}\\
A_3 = & 1 & 1 & (2) & 4 & 5 & 8 & 10 & 14 & 18 & 24 & 30 & (1-x^3)^{-2}\\
A_4 = & 1 & 1 & 2 & (4) & 6 & 10 & 14 & 22 & 30 & 44 & 60 & (1-x^4)^{-4}\\
A_5 = & 1 & 1 & 2 & 3 & (9) & 14 & 25 & 41 & 64 & 96 & 144 & (1-x^5)^{-9}\\
A_6 = & 1 & 1 & 2 & 3 & 5 & (20) & 29 & 56 & 95 & 153 & 247 & (1-x^6)^{-20}\\
A_7 = & 1 & 1 & 2 & 3 & 5 & 9 & (48) & 75 & 142 & 247 & 410 & (1-x^7)^{-48}\\
A_8 = & 1 & 1 & 2 & 3 & 5 & 9 & 20 & (115) & 181 & 348 & 624 & (1-x^8)^{-115}\\
A_9 = & 1 & 1 & 2 & 3 & 5 & 9 & 20 & 48 & (286) & 460 & 884 & (1-x^9)^{-286}\\
A_{10} = & 1 & 1 & 2 & 3 & 5 & 9 & 20 & 48 & 115 & (719) & 1842 & (1-x^{10})^{-719}\\
\hline
A_n = & 1 & 1 & 2 & 3 & 5 & 9 & 20 & 48 & 115 & 286 & 719 & \\
\text{for } n = & 1 & 2 & 3 & 4 & 5 & 6 & 7 & 8 & 9 & 10 & 11 & \\
\end{array}
\]

I have had occasion, for another purpose, to consider the question of finding the
number of trees with a given number of free branches, bifurcations at least. Thus,
when the number of free branches is three, the trees of the form in question are
those in the annexed figure, and the number is therefore two. It is not difficult to
see that we have in this case ($R_r$ being the number of such trees with $r$ free branches),
\[
(1 - x)^{-1}(1 - x^2)^{-R_2}(1 - x^3)^{-R_3}\cdots = 1 + x + 2R_2 x^2 + 2R_3 x^3 + 2R_4 x^4 + \&\textrm{c.};
\]
and a like process of development gives:
\[
\begin{array}{c|cccccc}
R_n = & 1 & 2 & 5 & 12 & 33 & 90\\
\text{for } n = & 2 & 3 & 4 & 5 & 6 & 7
\end{array}
\]

I may mention, in conclusion, that I was led to the consideration of the foregoing
theory of trees by Professor Sylvester's researches on the change of the independent
variables in the differential calculus.

\bigskip

\noindent\textit{2, Stone Buildings, January 2, 1856.}

\newpage


\section*{204. On a Problem in the Partition of Numbers}
\addcontentsline{toc}{section}{204. On a Problem in the Partition of Numbers}

\begin{center}
[From the \textit{Philosophical Magazine}, vol.~XIII.~(1857), pp.~245--248.]
\end{center}

\bigskip

I\textsc{t} is required to find the number of partitions into a given number of parts,
such that the first part is unity, and that no part is greater than twice the preceding
part.

Commencing to form the partitions in question, those are
\[
\begin{array}{l}
1\ 1\ 1\ 1\ 1\ 1\ 1\ 1\ 1\ \&\textrm{c}.;\\
1\ 2;\ 1\ 1\ 2;\ 1\ 1\ 1\ 2;\ 1\ 1\ 1\ 1\ 2;\ \&\textrm{c}.;\\
1\ 1\ 2\ 3;\ 1\ 1\ 2\ 4;\ \&\textrm{c}.;
\end{array}
\]
and if we were to proceed to the 4-partitions, each 3-partition ending in 1 would give
rise to two such partitions; each 3-partition ending in 2 to four such partitions; each
3-partition ending in 3 to six partitions; each 3-partition ending in 4 to eight
such partitions. We form in this manner the Table:
\[
\begin{array}{c|*{14}{c}|c}
\text{Number of} & 1 & 2 & 3 & 4 & 5 & 6 & 7 & 8 & 9 & 10 & 11 & 12 & 13 & 14 & \text{Totals}\\
\text{partitions} & & & & & & & & & & & & & & &\\
\hline
1\text{-partitions} & 1 & & & & & & & & & & & & & & 1\\
2\text{-partitions} & & 1 & 1 & & & & & & & & & & & & 2\\
3\text{-partitions} & & & 1 & 2 & 1 & 1 & & & & & & & & & 5\\
4\text{-partitions} & & & & 1 & 2 & 3 & 3 & 3 & 2 & 2 & 1 & 1 & & & 18\\
5\text{-partitions} & & & & & 1 & 2 & 3 & 5 & 6 & 8 & 10 & 11 & 10 & 11 & 100
\end{array}
\]
and we are thus led to the series
\[
1,\ 2,\ 5,\ 18,\ \&\textrm{c.},
\]
where, considering it as the first term of each series, the first differences of any series
are the terms twice repeated of the next preceding series: thus the differences of the
fourth series are 1, 1, 2, 3, 4, 6, 6. It is moreover clear that the first half of
each series is precisely the series which immediately precedes it; we need, in fact,
only consider a single infinite series 1, 2, 5, 6, \&c. It is to be remarked, moreover,
that in the column of totals, the total of any term is precisely the first number in
the next succeeding line.

Consider in general a series $A$, $B$, $C$, $D$, $E$, \&c., and a series $A'$, $B'$, $C'$, $D'$, $E'$,
\&c. derived from it as follows:
\begin{align*}
A' &= 2A,\\
B' &= 2B,\\
C' &= 2A + 2B,\\
D' &= 2A + 2B + 2C,\\
E' &= 2A + 2B + 2C + 2D,\\
F' &= 2A + 2B + 2C + 2D + 2E,\\
&\quad\&\textrm{c.};
\end{align*}
viz. the first differences of the series 0, $A'$, $B'$, $C'$, $D'$, $E'$, \&c. are $A$, $A$, $B$, $B$, $C$, $C$, $D$, $D$,
\&c. Multiplying by 1, $a$, $a^2$, \&c. and adding, we have
\[
A' + B' a + C' a^2 + D' a^3 + E' a^4 + \&\textrm{c.} = 2(A + Ba + Ca^2 + \&\textrm{c.})
\]
and if we form in a similar manner $A''$, $B''$, $C''$, $D''$, \&c. from $A'$, $B'$, $C'$, $D'$, and
so on, we have
\[
A'' + B'' a + C'' a^2 + \&\textrm{c.} = \frac{1 + a}{1 - a}(A + Ba + Ca^2 + \&\textrm{c.}),
\]
\[
A''' + B''' a + C''' a^2 + \&\textrm{c.} = \frac{1 + a}{1 - a}\cdot\frac{1 + a^2}{1 - a^2}(A + Ba + Ca^2 + \&\textrm{c.});
\]
and so on. While $A = 1$, and suppose that the process is repeated an indefinite
number of times, we have
\[
1 + 2a + 5a^2 + 2a^3 + \&\textrm{c.} = \frac{1 + a}{1 - a}\cdot\frac{1 + a^2}{1 - a^2}\cdot\frac{1 + a^3}{1 - a^3}\cdot\frac{1 + a^4}{1 - a^4}\ \&\textrm{c.}
\]
and the coefficients $1$, $2$, $5$, $18$, \&c. are precisely simple algorithms for the calculation of the
coefficients of $a$, $a^2$, \&c. which is the third notation of the following theorem:
\[
1 + 2a + 5a^2 + 18a^3 + 100a^4 + \&\textrm{c.} = \frac{1 + a}{(1 + a)^2 (1 + a^2)} - x^{42}(1 + a^2)^{-43} - x^{44}(1 + a^3)^{-44}\&\textrm{c.}.
\]
which gives also to the following very simple algorithm for finding the number of
coefficients of $a$, $a^2$, \&c. (in question are positive numbers as is the case);
\[
1 + 2a + 6a^2 + 2a^3 + \&\textrm{c.}\ ;
\]
this, on multiplying by 1, $a$, $a^2$, \&c. (in question are positive numbers and the doubling
$3a^2$, $4a^3$, \&c. and consequently, the operation $-3a^2(1 + a^2)^{-3} - 4a^3(1 + a^4)^{-3}\&\textrm{c}$. (we may write the
operation as a series, $1 + 2a + 5a^2 + 6a^3 + 6a^4 + \&\textrm{c}$.), then this product is the same as the
expression which represents the number of partitions of $n$, viz. as in the part question.

The proceeding expression for $1 + 2a + 5a^2 + 6a^3 + \&\textrm{c}$. as part greater than twice
the preceding part is equal to the number of partitions of $n + 1$ into parts not
greater than the previous one, may by itself become to the following:
\[
1 + 2a + 5a^2 + 2a^3 + 4a^4 + \&\textrm{c}.
\]
For example, the partitions of 5, 1, 2; are with the parts 1, 1, 7; 2, 3 = 1; the
number of 3-partitions is 4; and by the second part of the theorem, the number of
4-partitions is
\[
3(1 + 2 + 4 + 6) = 36.
\]
the number of which are 1, 2, 4. Hence, by the first part of the theorem, the
number of 5-partitions is 8, and by the second part of the theorem, the number of
3-partitions is 5.

\bigskip

\noindent\textit{2, Stone Buildings, March 17, 1857.}

\newpage


\section*{205. Note on the Summation of a Certain Factorial Expression}
\addcontentsline{toc}{section}{205. Note on the Summation of a Certain Factorial Expression}

\begin{center}
[From the \textit{Philosophical Magazine}, vol.~XIII.~(1857), pp.~419--423.]
\end{center}

\bigskip

M\textsc{r} K\textsc{irkman} some months ago communicated to me a formula for the double
summation of a factorial expression, to which formula he had been led by his researches
on the partition of polygons. The formula is a slightly altered form is as follows: viz.
\[
\sum_{m=0}^{m=n}\frac{(\rho + s + 2)^m(s)^m}{(\rho)^m} \cdot \frac{(\rho + s - 1 - 2)^{m-1}}{(s-1)^{m-1}} = \frac{2k}{(k+1)^{(k+1)}}(s)^{(k+1)}
\]
where the summation extending from $z = 0$ to $z = n - 1$. In the particular case where $k = 0$,
this is a formula well known. In the same case the series except those in which
$\rho = s$ vanish; and putting therefore $k = 0$, and making a slight change in
the form of the right-hand side, the formula becomes
\[
\sum_{m=0}^{m=n}\frac{(2\rho + 2k)^m}{\rho^m}\cdot \frac{(2\rho + 1)^{(\rho-1)}}{(2\rho+1)^{(\rho-1)}} = \frac{(2\rho + 1)^{(\rho-1)}}{(\rho-1)^{(\rho-1)}}
\]
the summation extending from $z = 0$ to $z = n - 1$.

\bigskip

We have, in the notation of Gauss, $[m]^z = m(m - 1)\ldots 2.1 = \Pi m$, and a factorial
$[m]^z$ is expressed in terms of the function $\Pi$ by the formula $[m]^z = \Pi m \div \Pi(m - n)$.
Write also
\[
\Pi(m - 1) = \Pi(m) \div [m].
\]
we have
\[
\Pi m = m \Pi(m - 1) = m \Pi(m - 1),
\]
\[
\Pi(2m + 1) = 2m \Pi(2m - 1).
\]

\bigskip

\noindent\textit{2, Stone Buildings, January 5, 1857.} \quad\textit{[continues on the next page set]}

\bigskip

\noindent\textit{[End of pp.~201--250 chunk.]}

\end{document}
